\documentclass[12pt]{article}
\usepackage[letterpaper, margin=1.5cm]{geometry}
\usepackage{amsmath}
\usepackage{amssymb}
\usepackage{empheq}
\setlength\parindent{0pt}

\begin{document}


Von Mises se define mediante la siguiente ecuación:
\[ f(\pmb{\sigma}) = \sqrt{\frac{3}{2}\pmb{S}:\pmb{S}}\]
\[ \pmb{S} = \pmb{\sigma} - \frac{1}{3}\text{tr}(\pmb{\sigma})\pmb{I} \]
Para calcular la primera derivada de la función, se utiliza la regla de la cadena:
\[ \frac{\partial f}{\partial \pmb{\sigma}} = \frac{\partial f}{\partial \pmb{S}} : \frac{\partial \pmb{S}}{\partial \pmb{\sigma}} \]
Donde:
\[ \frac{\partial f}{\partial \pmb{S}} = \frac{3}{2}\frac{\pmb{S}}{\sqrt{\pmb{S}:\pmb{S}}} = \frac{3}{2}\frac{\pmb{S}}{f} \]
\[ \frac{\partial \pmb{S}}{\partial \pmb{\sigma}} = \mathbb{I} - \frac{1}{3}\pmb{I} \otimes \pmb{I} \]

Se calcula la derivada completa:
\[ \frac{\partial f}{\partial \pmb{\sigma}} = \frac{3}{2}\frac{\pmb{S}}{f} : \left(\mathbb{I} - \frac{1}{3}\pmb{I} \otimes \pmb{I}\right) = \frac{3}{2f}\left(\pmb{S}:\mathbb{I} - \frac{1}{3}\pmb{S}:(\pmb{I} \otimes \pmb{I})\right) = \frac{3}{2f}\left(\pmb{S} - \frac{1}{3}(\pmb{S}:\pmb{I}) \otimes \pmb{I}\right)\]
Como \(\pmb{S}\) es un tensor desviador, su traza es cero, por lo que el segundo término desaparece, quedando:
\[ \frac{\partial f}{\partial \pmb{\sigma}} = \frac{3}{2f}\pmb{S} \]

Para la segunda derivada, se aplica nuevamente la regla de la cadena:
\[ \frac{\partial^2 f}{\partial \pmb{\sigma} \partial \pmb{\sigma}} = \frac{\partial}{\partial \pmb{\sigma}}\left(\frac{\partial f}{\partial \pmb{\sigma}}\right) = \frac{\partial}{\partial \pmb{\sigma}}\left(\frac{3}{2f}\pmb{S}\right) = \frac{\partial}{\partial \pmb{S}}\left(\frac{3}{2f}\pmb{S}\right): \frac{\partial\pmb{S}}{\partial \pmb{\sigma}} \]

Calculando la primera parte:
\[ \frac{\partial}{\partial \pmb{S}}\left(\frac{3}{2f}\pmb{S}\right) = \frac{3}{2}\left(\frac{\partial}{\partial \pmb{S}}\left(\frac{1}{f}\right)\pmb{S} + \frac{1}{f}\frac{\partial \pmb{S}}{\partial \pmb{S}}\right) \]
\[ \frac{\partial}{\partial \pmb{S}}\left(\frac{1}{f}\right) = -\frac{1}{f^2}\frac{\partial f}{\partial \pmb{S}} = -\frac{1}{f^2}\cdot \frac{3}{2}\frac{\pmb{S}}{f} = -\frac{3}{2f^3}\pmb{S} \]
Por lo tanto:
\[ \frac{\partial}{\partial \pmb{S}}\left(\frac{3}{2f}\pmb{S}\right) = \frac{3}{2}\left(-\frac{3}{2f^3}\pmb{S} \otimes \pmb{S} + \frac{1}{f}\mathbb{I}\right) = \frac{3}{2f}\mathbb{I} - \frac{9}{4f^3}\pmb{S} \otimes \pmb{S} \]
Ahora se multiplica por la segunda parte:
\[ \frac{\partial^2 f}{\partial \pmb{\sigma} \partial \pmb{\sigma}} = \left(\frac{3}{2f}\mathbb{I} - \frac{9}{4f^3}\pmb{S} \otimes \pmb{S}\right) : \left(\mathbb{I} - \frac{1}{3}\pmb{I} \otimes \pmb{I}\right) \]
Realizando la multiplicación:
\[ \frac{\partial^2 f}{\partial \pmb{\sigma} \partial \pmb{\sigma}} = \frac{3}{2f}\left(\mathbb{I} - \frac{1}{3}\pmb{I} \otimes \pmb{I}\right) - \frac{9}{4f^3}\left(\pmb{S} \otimes \pmb{S} - \frac{1}{3}(\pmb{S}:\pmb{I})\left(\pmb{S} \otimes \pmb{I}\right)\right) \]
Nuevamente, como \(\pmb{S}\) es un tensor desviador, su traza es cero, por lo que el segundo término dentro del paréntesis desaparece, quedando:
\[ \frac{\partial^2 f}{\partial \pmb{\sigma} \partial \pmb{\sigma}} = \frac{3}{2f}\left(\mathbb{I} - \frac{1}{3}\pmb{I} \otimes \pmb{I}\right) - \frac{9}{4f^3}\left(\pmb{S} \otimes \pmb{S}\right) \]
\end{document}
